\section{The Crystal: Coxeter Groups, Schl\"afli Symbols, and the Dodecagrid}

The previous section showed that the crystal of experience must be a regular tiling of hyperbolic space.
This section names that tiling exactly and shows how it is generated. The natural question is where an
infinite, perfectly regular crystal comes from. The answer is not that someone laid down cell after cell by
hand. The whole tiling is generated automatically, from one tile, by mirrors. This is the theory of Coxeter
groups. We give the reflection picture, the Schl\"afli notation that names the crystal in a few integers,
the curvature rule with a worked side-by-side calculation, the two crystals the theory works with, and the
argument that fixes the dimension at three.

\subsection{Coxeter groups as reflection groups}

Look into a kaleidoscope. A few scraps of colored glass sit between mirrors, and you see a vast, perfectly
symmetric pattern. The mirrors did all the work. Each mirror reflects the scrap to the other side, those
reflections have their own reflections, and the copies multiply until the whole field is full. A regular
tiling is made the same way. Take one tile, stand a mirror along each face, and each mirror reflects the
tile to a copy on the far side. Reflections of reflections fill the space with identical tiles. You never
placed a single one after the first.

Fix one tile, the fundamental chamber, and put a mirror on each of its faces. In three dimensions a chamber
is a simplex with four faces, so there are four mirrors, with reflections $s_0, s_1, s_2, s_3$. Each $s_i$
flips the chamber across face $i$ onto a fresh copy. A reflection done twice is the identity, so every
generator satisfies $s_i s_i = e$, and each generator is its own inverse. Composing reflections lands you on
some copy of the chamber, and every copy is reached by exactly such a sequence. The set of all these
sequences, under composition, is the Coxeter group. The tiling is what you see when you apply every group
element to one chamber.

One more family of relations carries the geometry. Two mirrors that share an edge meet at a dihedral angle
of $\pi / m$. Reflecting back and forth between them is a rotation, and after $m$ round trips it closes back
to the identity, so $(s_i s_j)^m = e$. The integer $m$, written $m_{ij}$, is the only thing chosen for that
pair of mirrors. When two mirrors are perpendicular, $m_{ij} = 2$ and the reflections commute.

The full presentation is generators $s_0, \dots, s_n$ with relations $(s_i s_j)^{m_{ij}} = e$, where
$m_{ii} = 1$ and $m_{ij} = m_{ji} \ge 2$ for $i \ne j$. A handful of integers fixes the angles, the angles
fix the mirrors, the mirrors generate the tiling, and the tiling fixes the curvature.

The cleanest fact in the subject is that one chamber equals one group element. The fundamental chamber is
the identity $e$, the chamber $s_i e$ is the one across mirror $i$, and so on. Every copy of the chamber
matches exactly one group element, and the chamber is a fundamental domain. To list the chambers is to list
the group elements, and to find a chamber's neighbors is to multiply by generators. Geometry becomes word
combinatorics, and the crystal can be enumerated with no coordinates at all, by reducing words to a
canonical ShortLex normal form (P85, P87).

\subsection{The $[5,3,4]$ group}

Assemble the pieces for the crystal of the theory. The Coxeter group $[5,3,4]$ has four generators and the
relations
\begin{align}
s_i s_i &= e && \text{for every } i, \\
(s_0 s_1)^5 &= e && \text{the pentagon angle}, \\
(s_1 s_2)^3 &= e && \text{the triangle angle}, \\
(s_2 s_3)^4 &= e && \text{four cells around an edge}, \\
(s_i s_j)^2 &= e && \text{perpendicular otherwise}.
\end{align}
The Gram matrix of the mirror normals has a signature with exactly one negative eigenvalue. That one
negative direction is the timelike axis, and the geometry is hyperbolic. The cell of the crystal is the
dodecahedron $\{5,3\}$, whose symmetry group is the icosahedral group $H_3$ of order $120$. The word engine
enumerates this cell stabilizer to exactly $120$ chambers, with no approximation. Around each edge of the
crystal four dodecahedra meet, the $r = 4$ in the symbol, and each cell has exactly $12$ neighbors, one
across each of its twelve pentagonal faces. Both the order $120$ and the neighbor count $12$ are confirmed
exactly (P85, P87).

The same one machine yields every other regular tessellation by changing the symbol alone (P47), so the
theory does not posit the crystal as a bespoke object. It posits a universal generator and a tiny integer
input. The crystal is the output. The general study of automata and computation on such hyperbolic tilings
is developed by Margenstern \cite{margenstern2013}.

\subsection{Schl\"afli symbols and the curvature rule}

A Schl\"afli symbol is a short list of whole numbers that names a regular shape, tiling, or honeycomb
completely. The symbol grows one number at a time, and each number says how the pieces from the level below
fit together at the next level up. One number, $\{p\}$, names a regular $p$-gon. Two numbers, $\{p,q\}$,
name $q$ copies of the $p$-gon around every vertex, which is either a polyhedron or a tiling. Three numbers,
$\{p,q,r\}$, name a honeycomb whose cells are the solids $\{p,q\}$, with $r$ of those cells packed around
every edge.

For a two-number symbol the curvature follows from a single check. Compute $1/p + 1/q$ and compare it to
$1/2$. If the sum is greater than $1/2$ the tiling closes into a ball and is spherical. If it equals $1/2$
the tiling lies flat and is Euclidean. If it is less than $1/2$ the tiling ruffles into a saddle and is
hyperbolic. The same fact in product form is to compare $(p-2)(q-2)$ to $4$, below being spherical, equal
flat, above hyperbolic.

Run it on $\{7,3\}$. One seventh plus one third is about $0.476$, below $0.5$, so it is hyperbolic, and
$(7-2)(3-2) = 5$, above $4$, agrees. Run it on $\{6,3\}$. One sixth plus one third is exactly $0.5$ and
$(6-2)(3-2) = 4$, so it is flat, right on the boundary. Seven is the smallest side count that bends the
plane at three faces to a corner.

\subsection{The worked calculation for three honeycombs}

For a three-number honeycomb the deciding question is how much angle the cells use up around an edge. Each
cell $\{p,q\}$ has a fixed dihedral angle $d$ along its edges, computed from
\begin{equation}
\sin\!\left(\frac{d}{2}\right) = \frac{\cos(\pi/q)}{\sin(\pi/p)}.
\end{equation}
Multiplying $d$ by $r$, the number of cells around an edge, and comparing the total to $360$ degrees decides
the geometry. A total under $360$ leaves a gap and closes up, so it is spherical. A total of exactly $360$
fills the angle and is Euclidean. A total over $360$ is too much angle for flat space, and the overflow is
absorbed by negative curvature, so it is hyperbolic. Three honeycombs show the dial walking from one regime
to the next.

For $\{4,3,4\}$, the cubic honeycomb, the cell $\{4,3\}$ is the cube. Here
$\sin(d/2) = \cos(\pi/3)/\sin(\pi/4) = 0.5/0.7071 = 0.7071$, so $d/2 = 45$ degrees and the dihedral is $90$
degrees. Four cubes around an edge give $4 \times 90 = 360$ exactly. The geometry is Euclidean. This is
ordinary flat three-dimensional space, tiled by cubes.

For $\{4,3,3\}$, the boundary of the tesseract, the cell is again the cube with a $90$ degree dihedral, but
only $r = 3$ cells meet at an edge. The total is $3 \times 90 = 270$ degrees, under $360$. The geometry is
spherical, and the honeycomb closes up into the finite four-cell-boundary of a hypercube.

For $\{5,3,4\}$, the dodecagrid, the cell is the dodecahedron. Its dihedral angle is about $116.57$ degrees,
and $r = 4$ cells meet at an edge, so the total is $4 \times 116.57 = 466.3$ degrees, well over $360$. There
is no way to fit four dodecahedra around an edge in flat space, since they would overlap. Curving the space
hyperbolically opens up exactly enough room for the fourth cell. So $\{5,3,4\}$ is forced into
three-dimensional hyperbolic space.

Changing a single number in the symbol walks the geometry from spherical at $270$ to flat at $360$ to
hyperbolic at $466$, and $\{5,3,4\}$ is the first dodecahedral honeycomb that overflows the edge. Once the
symbol is fixed there are no further parameters. The neighbor count, the symmetry group, the growth rate,
and the curvature all follow by computation.

\begin{table}[h]
\centering
\begin{tabular}{lllll}
\toprule
Honeycomb & Cell dihedral & Cells per edge & Total angle & Geometry \\
\midrule
$\{4,3,4\}$ & $90$ degrees & $4$ & $360$ degrees & Euclidean \\
$\{4,3,3\}$ & $90$ degrees & $3$ & $270$ degrees & Spherical \\
$\{5,3,4\}$ & $116.57$ degrees & $4$ & $466.3$ degrees & Hyperbolic \\
\bottomrule
\end{tabular}
\caption{The curvature dial for three honeycombs.}
\label{tab:curvature-dial}
\end{table}

\subsection{The two crystals}

The theory works with two crystals, one dimension apart. The first is a drawable study model, the second is
the real substrate. The heptagrid $\{7,3\}$ is the regular tiling of the hyperbolic plane by seven-sided
cells, three meeting at every corner. Each cell has exactly seven neighbors, one across each of its seven
edges, confirmed exactly in the same facet experiment that pins the dodecagrid at twelve (P85).

Drawn in the Poincare disk it is one of the cleanest pictures in mathematics, a central heptagon ringed by
more heptagons, each ring holding far more cells than the last, all appearing to shrink toward the rim. The
shrinking is an artifact of the drawing. In the true geometry every cell is the same size.

The heptagrid is the picture to hold when imagining the crystal of experience. Color each cell by its tone
and let two touching cells mean two feelings experiencing each other through their shared edge, and the
colored heptagrid is a faithful, literal picture of the model, just in two dimensions instead of three.

\begin{figure}[h]\centering\includegraphics[width=0.47\textwidth]{vibe-mesh-73.pdf}\caption{The heptagrid, the regular tiling of the hyperbolic plane by seven-sided cells, three to a corner, the simplest face of the crystal. Each cell is a vibe and each shared edge is a perceived relation. Cells appear to shrink toward the rim only because the Poincare disk compresses the infinite plane into a finite picture. In the true geometry every cell is the same size.}\label{fig:heptagrid}\end{figure}

The real substrate is one dimension up. It is the dodecagrid $\{5,3,4\}$, a packing of dodecahedra filling
curved three-dimensional space. The symbol reads as a chain. The $5$ means each face of a cell is a
pentagon, the $3$ means three pentagons meet at each corner of a cell, which assembles a dodecahedron, and
the $4$ means four dodecahedra meet around every edge. Each dodecahedral cell has exactly $12$ neighbors,
one across each of its twelve pentagonal faces, and this is the central number of the whole theory (P85,
P104).

You cannot fully picture this crystal, and that is fine. It is the exact three-dimensional analogue of the
heptagrid. Everything true of the heptagrid is true of the dodecagrid, with regular identical cells, no
preferred direction, exponential room, and the same kaleidoscope of reflections. The build is exact to tens
of millions of cells by an integer fingerprint method that avoids any floating-point drift (P104).

The heptagrid is the version you can see. The dodecagrid is the real substrate. Time is its growth, the
birth of new cells at the frontier, generation after generation, with each new layer the next moment.

\begin{figure}[h]\centering\includegraphics[width=\textwidth]{vibe-mesh-534.png}\caption{The dodecagrid $\{5,3,4\}$, the three-dimensional substrate, dodecahedral cells with four meeting around every edge, filling curved three-dimensional space. It is the heptagrid one dimension up, the crystal of which physical space is made. Time is its growth.}\label{fig:dodecagrid}\end{figure}

\subsection{Why three dimensions}

The dimension of space is not a dial in this theory. It is forced, and forced twice over, then narrowed
further to the specific crystal. The first squeeze is pure geometry. A crystal here is a regular tiling of
hyperbolic space with finite cells, a compact regular hyperbolic honeycomb. Such a honeycomb exists only
when both its cell and its vertex figure are themselves finite, which is a sharp constraint on the symbol.

Experiment P62 enumerates every such symbol by Gram-matrix signatures and finds a hard pattern. The
hyperbolic plane admits infinitely many. Three-dimensional hyperbolic space admits exactly four, one of
which is the dodecagrid $\{5,3,4\}$. Four-dimensional hyperbolic space admits exactly five. Five dimensions
and above admit none, ever. The window of nonzero dimensions is precisely two, three, and four. A
finite-celled crystal substrate simply cannot be built above four dimensions, so there is no room for hidden
spatial dimensions to hide.

\begin{table}[h]
\centering
\begin{tabular}{p{2.1cm}p{3.1cm}p{7.2cm}}
\toprule
Dimension & Compact regular hyperbolic honeycombs & Examples \\
\midrule
$2$ & infinitely many & all $\{p,q\}$ with $1/p + 1/q < 1/2$ \\
$3$ & four & $\{3,5,3\}$, $\{4,3,5\}$, $\{5,3,4\}$, $\{5,3,5\}$ \\
$4$ & five & $\{3,3,3,5\}$, $\{4,3,3,5\}$, $\{5,3,3,3\}$, $\{5,3,3,4\}$, $\{5,3,3,5\}$ \\
$5$ and above & none & none \\
\bottomrule
\end{tabular}
\caption{The dimension window. Compact regular hyperbolic crystals exist only in dimensions $2$, $3$, and
$4$, with the dodecagrid $\{5,3,4\}$ the chosen one.}
\label{tab:dimension-window}
\end{table}

The second squeeze is physics and picks one of the three survivors. For stable matter to exist, orbits must
close and retrace themselves, otherwise nothing stays bound over time. Whether an orbit closes depends on
how force falls off with distance, which depends on the dimension. Experiment P68 integrates these orbits
across the allowed window and measures the apsidal angle, the angle swept from one closest approach to the
next.

In two dimensions the orbit precesses and never quite closes, too loose to hold structure. In three
dimensions the inverse-square law gives an apsidal angle of exactly $180$ degrees, so the orbit retraces the
same ellipse, the world of stable atoms and solar systems. In four dimensions and above the orbit is
unstable, spiraling in or flying off at the slightest nudge.

Geometry allows two, three, or four. Physics allows only three. The intersection is a single number. A
confirming measurement grows a crystal from scratch and reads back its intrinsic dimension from connectivity
alone, returning values near the targets without bias (P69).

The two squeezes fix the dimension at three, but three-dimensional hyperbolic space still has four compact
crystals. The auto-selection result shows $\{5,3,4\}$ is derived rather than chosen from that menu (P86).
The ternary tone of each vibe makes the cell vertices three-valent, which is the middle entry $q = 3$. The
finite three-valent cells are exactly the tetrahedron $\{3,3\}$, the cube $\{4,3\}$, and the dodecahedron
$\{5,3\}$.

Eternal growth requires a genuinely hyperbolic crystal that does not close into a finite ball, and the
theory takes the tightest such crystal, the one with the fewest cells per edge. The tetrahedron reaches no
compact crystal at all. The cube needs $r = 5$. The dodecahedron reaches a compact hyperbolic crystal at
$r = 4$, the tightest closure across all three-valent candidates. The forced symbol is exactly $\{5,3,4\}$.

The golden ratio that builds the dodecahedron is not smuggled in. The five falls out of the requirement to
be the tightest eternal crystal with three-valent cells. The crystal at the root of reality is the unique
survivor of constraints the theory committed to for other reasons, and the dimension is a result rather than
an input.
